\section{Implementation}

\subsection{Adaptation Model}

\subsection{Grapheme-Phoneme Alignment}

Graphemes were matched to corresponding phonemes based on the
surrounding context of a each individual token. By default it is
a one-to-one relation, where one phoneme maps to one grapheme
according to shared vector indices, but often that is not the case
as phonetic structures are varied and highly dependent on surrounding
orthography or phonetics~\cite{Tayao_2004}. With consideration to this,
context-sensitive rules are made to produce an AlignedString. Type of
AlignedString consists of an array of two tuples, grapheme tokens taken 
from the English input word, and an array of IPA phoneme tokens returned from
the English G2P process. Each index consists of a singular grapheme
with its matched vector of IPA phonemes. For example, input string "hello"
has the IPA of \textipa{h@loU}, mapped as

\begin{enumerate}
    \item \( \text{h} \rightarrow \text{\textipa{h}}\)
    \item \( \text{e} \rightarrow \text{\textipa{@}}\)
    \item \( \text{l} \rightarrow \text{\textipa{l}}\)
    \item \( \text{l} \rightarrow None \)
    \item \( \text{o} \rightarrow \text{\textipa{oU}}\)
\end{enumerate}

Individual graphemes are mapped to their corresponding phonetic equivalent,
and context-sensitive alignments are gauged on the surrounding graphemes or
phonemes to more accurately map corresponding grapheme to phoneme pairs.
This process can be broken down into separate context-sensitive grapheme
and phoneme rules, where grapheme rules rely solely on the surrounding graphemes given an indexed token.
Phoneme rules are further broken down into non-consuming and consuming cases, where
non-consuming allows the current phoneme to be of value `None` without affecting
the indexing order while consuming cases allow multiple phonemes to merge
into a singular token represented by a vector.

Alignment was needed as a preliminary step towards the adaptation of English loanwords
into Philippine English due to the variability of sounds in relation to written orthography.

\subsection{Vowel Mapping}

